\section{本讲主线：用选择性计算控制成本}

本节建立两条主线：attention alternatives 用更便宜的状态/稀疏机制控制长上下文成本；MoE 用条件计算扩大参数量但控制每 token FLOPs。



\noindent\textbf{展开说明：}A T T E N T I O N A LT E R N AT I V E S A N D M I X T U R E S O F E X P E R T S



\begin{figure}[H]
\centering
\includegraphics[width=0.88\textwidth]{slide-001.jpg}
\caption{Slide 2：Attention alternatives 的动机；上下文变长后 full attention 成本快速增长。}
\figsource{CS336 Spring 2026 Lecture 4 官方幻灯片。}
\end{figure}

\noindent\textbf{展开说明：}Cost of attention rises with large context sizes… how do we control those costs? 读这页时应把问题拆成三种资源：训练时的二次 attention FLOPs、推理时随序列增长的 KV/state memory，以及跨设备搬运状态的通信。Linear/recurrent attention、local/global hybrid 和 sparse selection 分别压缩不同资源，因此不能只用一个“复杂度更低”概括。


\begin{figure}[H]
\centering
\includegraphics[width=0.88\textwidth]{slide-002.jpg}
\caption{Slide 3: The ‘basic’ toolkit}
\figsource{CS336 Spring 2026 Lecture 4 官方幻灯片。}
\end{figure}

\noindent\textbf{展开说明：}Combine local + global attention                     Systems engineering

\subsection{本章小结}
本节的共同线索是 selective computation：通过结构选择让 token 不必总是访问所有历史或所有参数。但选择性越强，越需要额外机制保证信息流、负载均衡和训练稳定。

